\documentclass[11pt]{article}
\usepackage{series}   % house style
\usepackage{amsmath,amssymb,amsthm}
\usepackage{geometry}
\usepackage{hyperref}
\usepackage{enumitem}

\geometry{margin=1in}

% ---------- metadata ----------
\title{The Modal Triplet Theory Program A1: Coherent Kinematics}
\author{Peter Nero}
\date{January 2026}

\begin{document}

\maketitle

\begin{abstract}
We develop a kinematic framework for Modal Triplet Theory (MTT) that does not
presuppose spacetime, geometry, or equations of motion. Kinematics is defined
structurally, in terms of persistence of coherent structures across overlapping
admissible encodings in the encoding atlas. Motion is identified with admissible
continuation through chains of overlapping descriptions, while worldlines are
equivalence classes of such continuations under re-encoding.

We show that causal structure, horizon formation, and irreversibility arise as
inevitable consequences of finite admissibility and obstruction to global
description. Null and timelike distinctions are derived from admissible
continuation properties rather than from a metric. No dynamical laws are assumed;
the results apply to all realization classes compatible with the MTT structural
core.

This work provides the kinematic layer of the MTT program, bridging the abstract
classification of reduced descriptions and the downstream encoding-class papers
in which gravity, gauge structure, and quantization emerge as bookkeeping
responses to structural obstructions.
\end{abstract}

\tableofcontents
\newpage

% ============================================================
\section{Introduction and Scope}

The structural core of Modal Triplet Theory establishes when reduced
descriptions exist, how they relate on overlaps, and why no single global
description can exist. These results are deliberately pre-kinematic: they do
not assume spacetime, trajectories, or equations of motion. Nevertheless,
any theory that admits persistent reduced descriptions must address a basic
question: \emph{what does it mean for something to move, persist, or propagate
when no global kinematic arena exists?}

The purpose of the present paper is to answer this question at the structural
level. We develop a notion of kinematics that is intrinsic to the encoding
atlas of Modal Triplet Theory. In this framework, kinematics is not defined by
motion through a background space, but by persistence of coherent structure
across admissible overlaps of local descriptions.

Specifically, we show that:
\begin{itemize}
\item \emph{Position} is an encoding-relative notion defined by support of
coherent structure within a local description.
\item \emph{Motion} corresponds to admissible continuation through chains of
overlapping encodings.
\item \emph{Worldlines} are equivalence classes of admissible continuations
under re-encoding.
\item \emph{Horizons} and \emph{termination of motion} arise when admissible
continuation fails (nil obstructions).
\item \emph{Irreversibility} follows from the absence of global sections and the
terminal nature of nil.
\end{itemize}

Throughout this work, no metric, connection, or dynamical law is assumed. The
results are purely structural and apply uniformly across all realization
classes compatible with the MTT core. Geometry and dynamics, when they appear
in later papers, are shown to encode and stabilize the kinematic structures
identified here rather than to define them.

\begin{remark}[Relation to the series]
This paper depends only on the structural core of Modal Triplet Theory and the
encoding atlas it defines. It precedes the encoding-class papers on gravity,
gauge structure, and quantization, and provides the kinematic interpretation
required to understand those encodings as responses to structural obstructions
rather than as fundamental postulates.
\end{remark}

% ============================================================
% Next section will be:
% \section{Encoding-Relative Configuration and Position}
% ============================================================
\section{Encoding-Relative Configuration and Position}

We begin by defining the kinematic notions of configuration and position without
reference to a background space. All definitions are relative to admissible
encodings in the encoding atlas.

\subsection{Encoding-relative configuration}

Let $\mathcal E_\alpha = (A_\alpha, Z_\alpha, E_\alpha, \ldots)$ be an admissible
encoding on domain $A_\alpha \subset X$.

\begin{definition}[Encoding-relative configuration]
The \emph{configuration} of a coherent structure with respect to an encoding
$\mathcal E_\alpha$ is its image in the reduced variable space $Z_\alpha$ under
the encoding map $E_\alpha$.
\end{definition}

Configurations are therefore encoding-dependent. There is no notion of a
configuration independent of an encoding.

\begin{remark}
Different encodings may assign different configurations to the same coherent
structure. Overlap consistency ensures that these configurations are related by
admissible re-encoding on intersections of admissible domains.
\end{remark}

\subsection{Support of coherent structure}

The notion of position requires a notion of localization within an encoding.

\begin{definition}[Support]
Let $B \subset A_\alpha$ be the subset of the admissible domain on which a given
coherent structure is represented. The \emph{support} of the structure in the
encoding $\mathcal E_\alpha$ is the image
\[
\mathrm{supp}_\alpha := E_\alpha(B) \subset Z_\alpha.
\]
\end{definition}

Support captures where, within the encoding-relative configuration space, the
coherent structure is localized.

\subsection{Position as an encoding-relative diagnostic}

We now define position in purely kinematic terms.

\begin{definition}[Position]
A \emph{position} in encoding $\mathcal E_\alpha$ is any admissible diagnostic
functional
\[
\mathrm{pos}_\alpha : \mathcal P(Z_\alpha) \to \mathcal Q_\alpha
\]
that assigns to the support $\mathrm{supp}_\alpha$ a value in some diagnostic
space $\mathcal Q_\alpha$ (for example, a coordinate tuple, region label, or
coarse localization class).
\end{definition}

Position is therefore:
\begin{itemize}
\item encoding-relative,
\item diagnostic rather than fundamental,
\item defined only where an encoding is admissible.
\end{itemize}

\begin{remark}
No assumption is made that $\mathcal Q_\alpha$ carries a metric, topology, or
linear structure. Such structure may appear in particular realizations but is
not required at the kinematic level.
\end{remark}

\subsection{Consistency of position on overlaps}

On overlaps of admissible domains, positions defined in different encodings must
be mutually consistent.

\begin{definition}[Position compatibility]
Let $\mathcal E_\alpha$ and $\mathcal E_\beta$ be admissible encodings with
nonempty overlap $A_{\alpha\beta}$. Position diagnostics
$\mathrm{pos}_\alpha$ and $\mathrm{pos}_\beta$ are \emph{compatible} if there
exists an admissible re-encoding map $f_{\alpha\beta}$ such that
\[
\mathrm{pos}_\beta\bigl(E_\beta(B)\bigr)
\;\sim\;
\mathrm{pos}_\alpha\bigl(E_\alpha(B)\bigr)
\]
for all coherent supports $B \subset A_{\alpha\beta}$, where $\sim$ denotes
admissible equivalence.
\end{definition}

This condition ensures that different local notions of position do not contradict
one another on overlaps.

\subsection{Absence of global position}

Finite admissibility implies that no single global notion of position exists.

\begin{theorem}[No global position]
If the admissible domains do not admit a global encoding, then there exists no
globally defined position diagnostic consistent with all admissible encodings.
\end{theorem}

\begin{proof}
A globally defined position would require a global encoding whose restriction to
each admissible domain reproduces the local position diagnostics. This contradicts
finite admissibility.
\end{proof}

\begin{remark}
The absence of global position is not a deficiency of the theory but a structural
feature. It reflects the impossibility of global reduced description established
in the MTT core.
\end{remark}

\subsection{Preview: from position to motion}

Position as defined above is purely static and encoding-relative. In the next
section we introduce \emph{admissible continuation} across overlapping encodings
and define motion as persistence of coherent structure along such continuations.

% ============================================================
% Next section will be:
% \section{Admissible Continuation and Motion}
% ============================================================
\section{Admissible Continuation and Motion}

Having defined configuration and position relative to admissible encodings, we
now introduce the central kinematic notion of \emph{motion}. In Modal Triplet
Theory, motion is not defined as change of position in a background space, but as
persistence of coherent structure across chains of overlapping admissible
encodings.

\subsection{Admissible continuation}

Let $\mathcal E_\alpha = (A_\alpha, Z_\alpha, E_\alpha, \ldots)$ and
$\mathcal E_\beta = (A_\beta, Z_\beta, E_\beta, \ldots)$ be admissible encodings
with nonempty overlap $A_{\alpha\beta}$.

\begin{definition}[Admissible continuation]
A coherent structure represented in $\mathcal E_\alpha$ is said to admit an
\emph{admissible continuation} to $\mathcal E_\beta$ if there exists a subset
$B \subset A_{\alpha\beta}$ such that the representation in $Z_\alpha$ can be
re-encoded consistently into $Z_\beta$ via the admissible re-encoding map
$f_{\alpha\beta}$.
\end{definition}

Admissible continuation is therefore a relation between encodings, not a map on a
fixed configuration space.

\subsection{Chains of continuation}

Motion arises when admissible continuation exists along a sequence of overlaps.

\begin{definition}[Continuation chain]
A \emph{continuation chain} is a finite or infinite sequence of admissible
encodings
\[
\mathcal E_{\alpha_1}, \mathcal E_{\alpha_2}, \ldots
\]
such that admissible continuation exists from $\mathcal E_{\alpha_i}$ to
$\mathcal E_{\alpha_{i+1}}$ for all $i$.
\end{definition}

Continuation chains encode the possibility of persistent description without
requiring a global kinematic arena.

\subsection{Motion as persistence}

We now define motion itself.

\begin{definition}[Motion]
A coherent structure is said to be \emph{in motion} if it admits an admissible
continuation along a nontrivial continuation chain of encodings.
\end{definition}

Motion is thus:
\begin{itemize}
\item encoding-relative,
\item local to admissible domains,
\item defined by persistence of representation rather than by displacement.
\end{itemize}

\begin{remark}
A structure that admits continuation only within a single encoding but not
across overlaps is stationary in that encoding, but may fail to be stationary
globally due to the absence of global coherence.
\end{remark}

\subsection{Directionality and ordering}

Continuation chains naturally induce an ordering, but not a global time
parameter.

\begin{definition}[Kinematic ordering]
A \emph{kinematic ordering} is the partial order on admissible encodings induced
by admissible continuation: $\mathcal E_\alpha \prec \mathcal E_\beta$ if
admissible continuation exists from $\mathcal E_\alpha$ to $\mathcal E_\beta$.
\end{definition}

This ordering replaces external time in the kinematic description.

\begin{remark}
The kinematic ordering is generally not total. Branching, merging, and terminal
chains are possible, reflecting the absence of global coherence.
\end{remark}

\subsection{Reversibility and local symmetry}

Admissible continuation need not be symmetric.

\begin{definition}[Local reversibility]
A continuation from $\mathcal E_\alpha$ to $\mathcal E_\beta$ is \emph{locally
reversible} if admissible continuation also exists from $\mathcal E_\beta$ back
to $\mathcal E_\alpha$.
\end{definition}

\begin{remark}
Local reversibility may hold within restricted regions of the encoding atlas,
but fails globally when nil obstructions or selection fronts are encountered.
This asymmetry underlies the emergence of irreversibility.
\end{remark}

\subsection{Motion without background space}

We emphasize that no notion of distance, velocity, or trajectory in a background
space has been assumed.

\begin{remark}
Motion as defined here is compatible with many downstream realizations,
including geometric ones, but does not presuppose them. Any appearance of
spacetime motion in later encoding classes arises as an interpretation of
continuation structure, not as a primitive kinematic input.
\end{remark}

\subsection{Preview: worldlines}

Continuation chains are not unique: different chains may represent the same
persistent structure under re-encoding. In the next section we define
\emph{worldlines} as equivalence classes of admissible continuation chains.

% ============================================================
% Next section will be:
% \section{Worldlines as Equivalence Classes of Continuation}
% ============================================================
\section{Worldlines as Equivalence Classes of Continuation}

In the absence of a global kinematic arena, the notion of a worldline must be
reformulated. In Modal Triplet Theory, a worldline is not a curve in space or
spacetime, but an equivalence class of admissible continuation chains representing
the persistence of a coherent structure across overlapping descriptions.

\subsection{Non-uniqueness of continuation chains}

Given a coherent structure represented in an admissible encoding
$\mathcal E_\alpha$, there may exist multiple admissible continuation chains
emanating from $\mathcal E_\alpha$ that track the same structure across different
sequences of overlaps.

\begin{remark}
Distinct continuation chains may differ in:
\begin{itemize}
\item the choice of intermediate encodings;
\item the order in which overlaps are traversed;
\item the protocol or admissible re-encoding used on overlaps.
\end{itemize}
Despite these differences, the chains may represent the same persistent coherent
structure.
\end{remark}

This non-uniqueness reflects the absence of a preferred global parametrization.

\subsection{Equivalence of continuation chains}

We now formalize when two continuation chains represent the same worldline.

\begin{definition}[Worldline equivalence]
Two admissible continuation chains
\[
\{\mathcal E_{\alpha_1}, \mathcal E_{\alpha_2}, \ldots\}
\quad \text{and} \quad
\{\mathcal E_{\beta_1}, \mathcal E_{\beta_2}, \ldots\}
\]
are said to be \emph{equivalent} if, for every finite segment of one chain, there
exists a finite segment of the other such that the corresponding representations
are related by admissible re-encoding on overlapping domains.
\end{definition}

Equivalence is therefore defined by mutual representability rather than by
pointwise identification.

\subsection{Definition of worldline}

\begin{definition}[Worldline]
A \emph{worldline} is an equivalence class of admissible continuation chains under
the equivalence relation defined above.
\end{definition}

Worldlines encode kinematic identity without reference to background space,
time, or parametrization.

\subsection{Worldlines and the atlas nerve}

Continuation chains correspond naturally to paths in the atlas nerve.

\begin{remark}
A continuation chain defines a path in the $1$-skeleton of the atlas nerve. Worldline
equivalence identifies paths that differ by homotopy through admissible
re-encoding. In this sense, worldlines correspond to homotopy classes of
admissible paths in the atlas nerve.
\end{remark}

This interpretation makes clear that worldlines are global objects defined by
the overlap structure of admissible descriptions.

\subsection{Branching and merging}

Worldlines need not be unique or isolated.

\begin{definition}[Branching and merging]
A worldline \emph{branches} if there exist inequivalent continuation classes that
coincide up to some encoding and then diverge. A worldline \emph{merges} if
distinct classes become equivalent after a common continuation.
\end{definition}

\begin{remark}
Branching and merging reflect the partial nature of kinematic ordering and the
absence of a total time parameter. They are generic features in the presence of
finite admissibility.
\end{remark}

\subsection{Terminal worldlines}

Worldlines may terminate.

\begin{definition}[Terminal worldline]
A worldline is \emph{terminal} if it admits no admissible continuation beyond a
finite chain. Equivalently, its continuation encounters a nil obstruction.
\end{definition}

Terminal worldlines encode horizons, collapse, and termination of motion without
requiring singular behavior of the underlying system.

\subsection{Irreversibility of worldlines}

The absence of global coherence implies that worldline equivalence classes are
generally not reversible.

\begin{theorem}[Irreversibility of worldlines]
If a worldline is terminal or passes through a region of nil obstruction, then no
admissible continuation exists that retraces the worldline beyond that point.
\end{theorem}

\begin{proof}
By definition of nil obstruction, no admissible encoding exists beyond the
terminal point. Therefore no continuation chain can be extended through or
beyond that region, and reversal is impossible.
\end{proof}

\begin{remark}
Irreversibility here is kinematic rather than dynamical. It arises from the
structure of admissible descriptions, not from time-asymmetric equations of
motion.
\end{remark}

\subsection{Preview: causal structure}

Worldlines provide the foundation for defining causal relations among coherent
structures. In the next section we show how admissible continuation induces a
causal partial order and how null and timelike distinctions emerge from kinematic
constraints.

% ============================================================
% Next section will be:
% \section{Causal Structure and Kinematic Classification}
% ============================================================
\section{Causal Structure and Kinematic Classification}

Having defined motion as admissible continuation and worldlines as equivalence
classes of continuation chains, we now extract causal structure directly from
the kinematic ordering induced by admissibility. No background time, metric, or
dynamical law is assumed.

\subsection{Causal precedence}

Admissible continuation induces a partial causal order.

\begin{definition}[Causal precedence]
Given two admissible encodings $\mathcal E_\alpha$ and $\mathcal E_\beta$, we say
that $\mathcal E_\alpha$ \emph{causally precedes} $\mathcal E_\beta$, written
\[
\mathcal E_\alpha \prec \mathcal E_\beta,
\]
if there exists an admissible continuation chain from $\mathcal E_\alpha$ to
$\mathcal E_\beta$.
\end{definition}

This relation is reflexive and transitive but not necessarily antisymmetric,
reflecting the absence of global coherence.

\begin{remark}
Causal precedence is defined on encodings, not on an external time parameter.
It is therefore intrinsic to the structure of admissible descriptions.
\end{remark}

\subsection{Causal relations between worldlines}

Causal relations extend naturally to worldlines.

\begin{definition}[Worldline causality]
A worldline $W_1$ causally precedes a worldline $W_2$ if there exist
representatives $\mathcal E_\alpha \in W_1$ and $\mathcal E_\beta \in W_2$ such
that $\mathcal E_\alpha \prec \mathcal E_\beta$.
\end{definition}

This defines a partial order on worldlines modulo equivalence.

\subsection{Null and timelike continuation}

We now distinguish different classes of admissible continuation.

\begin{definition}[Null continuation]
An admissible continuation from $\mathcal E_\alpha$ to $\mathcal E_\beta$ is
\emph{null} if any admissible continuation chain between them is minimal, in the
sense that no intermediate admissible encoding can be inserted without losing
admissibility.
\end{definition}

\begin{definition}[Timelike continuation]
An admissible continuation from $\mathcal E_\alpha$ to $\mathcal E_\beta$ is
\emph{timelike} if there exists at least one admissible continuation chain
between them containing intermediate encodings.
\end{definition}

\begin{remark}
Null and timelike classifications are kinematic and encoding-relative. They do
not presuppose a metric or causal cone, but arise from the structure of
admissible continuation itself.
\end{remark}

\subsection{Causal cones}

The distinction between null and timelike continuation induces a cone-like
structure.

\begin{definition}[Causal cone]
The \emph{causal cone} of an encoding $\mathcal E_\alpha$ is the set of encodings
$\mathcal E_\beta$ such that $\mathcal E_\alpha \prec \mathcal E_\beta$.
\end{definition}

Null continuations form the boundary of the causal cone, while timelike
continuations form its interior.

\begin{remark}
This cone structure is purely order-theoretic and does not rely on any notion of
distance or angle. Its resemblance to light cones is a derived feature that
appears in suitable realizations.
\end{remark}

\subsection{Horizons as kinematic boundaries}

Finite admissibility implies that causal cones may be bounded.

\begin{definition}[Kinematic horizon]
A \emph{kinematic horizon} is a boundary in the causal cone of a worldline beyond
which no admissible continuation exists.
\end{definition}

\begin{remark}
Kinematic horizons arise when continuation chains terminate in nil
obstructions. They are not assumed; they follow directly from finite
admissibility.
\end{remark}

\subsection{Global causal failure}

The absence of a global encoding implies that causal structure is itself only
locally defined.

\begin{theorem}[No global causal ordering]
There exists no globally defined total causal order on all worldlines consistent
with admissible continuation.
\end{theorem}

\begin{proof}
A global total order would require a global encoding in which all admissible
continuations could be embedded. This contradicts finite admissibility.
\end{proof}

\subsection{Irreversibility revisited}

Causal structure provides an additional perspective on irreversibility.

\begin{remark}
Irreversibility arises because causal precedence is not invertible across
kinematic horizons and nil obstructions. This is a structural property of the
encoding atlas rather than a dynamical asymmetry.
\end{remark}

\subsection{Preview: relation to geometry}

The kinematic causal structure derived here will later be shown to require
additional bookkeeping to remain consistent under refinement. In particular,
the stabilization of causal cones and null boundaries motivates the emergence of
geometric and gravitational encoding structures in subsequent papers.

% ============================================================
% Next section will be:
% \section{Horizons, Selection, and Termination}
% ============================================================
\section{Horizons, Selection, and Termination}

In this section we analyze the kinematic consequences of finite admissibility.
We show how horizons, selection events, and termination of motion arise
structurally from the encoding atlas, without invoking dynamics, measurement
postulates, or external time.

\subsection{Finite admissibility and kinematic reach}

Finite admissibility implies that admissible continuation cannot persist
indefinitely along all directions.

\begin{definition}[Kinematic reach]
The \emph{kinematic reach} of an encoding $\mathcal E_\alpha$ is the set of
encodings $\mathcal E_\beta$ such that $\mathcal E_\alpha \prec \mathcal E_\beta$
via admissible continuation.
\end{definition}

Kinematic reach is generally bounded due to the presence of nil obstructions.

\subsection{Horizons as boundaries of continuation}

We now formalize horizons purely kinematically.

\begin{definition}[Kinematic horizon]
A \emph{kinematic horizon} relative to an encoding $\mathcal E_\alpha$ is a
boundary point of its kinematic reach such that:
\begin{enumerate}[label=(\roman*)]
\item admissible continuation exists up to the boundary;
\item no admissible continuation exists beyond the boundary.
\end{enumerate}
\end{definition}

\begin{remark}
Kinematic horizons are not defined by distances or metrics. They are defined by
the termination of admissible continuation and therefore exist even in the
absence of geometric structure.
\end{remark}

\subsection{Selection fronts}

Horizon formation is typically associated with abrupt changes in admissibility.

\begin{definition}[Selection front]
A \emph{selection front} is a codimension-one subset of the encoding atlas (or of
its nerve) across which the admissibility of continuation changes discontinuously,
forcing a transition between distinct admissible continuations or between
admissible continuation and termination.
\end{definition}

Selection fronts are the loci at which the system must ``choose'' among remaining
admissible encodings.

\subsection{Selection events}

We now define selection in kinematic terms.

\begin{definition}[Selection event]
A \emph{selection event} occurs when a continuation chain reaches a selection
front and admissible continuation persists only along a strict subset of the
previously available directions.
\end{definition}

\begin{remark}
Selection does not introduce randomness or collapse by assumption. It reflects
the structural loss of admissible continuation options.
\end{remark}

\subsection{Termination and nil revisited}

When no admissible continuation remains, termination occurs.

\begin{definition}[Termination]
A continuation chain \emph{terminates} if it encounters a nil obstruction, i.e.,
a region with no admissible encoding.
\end{definition}

Termination is therefore the kinematic manifestation of nil.

\subsection{Irreversibility from selection and termination}

Selection and termination together imply irreversibility.

\begin{theorem}[Kinematic irreversibility]
Once a continuation chain crosses a selection front or terminates at a nil
obstruction, there exists no admissible continuation that retraces the chain
beyond that point.
\end{theorem}

\begin{proof}
By definition of selection fronts, admissible continuation options are strictly
reduced. By definition of nil, no continuation exists. In either case, the set of
admissible continuations is not symmetric under reversal.
\end{proof}

\begin{remark}
Irreversibility here is kinematic and structural. It does not rely on entropy,
coarse-graining, or time-asymmetric dynamics.
\end{remark}

\subsection{Records as persistent kinematic traces}

Selection and termination give rise to records.

\begin{definition}[Record]
A \emph{record} is a persistent feature of a continuation chain that remains
invariant under all subsequent admissible continuations.
\end{definition}

\begin{remark}
Records arise when admissible continuation options narrow. They encode past
selection events structurally, not dynamically.
\end{remark}

\subsection{Preview: predictive limits}

Selection fronts and horizons place intrinsic limits on prediction.

\begin{remark}
The existence of selection fronts implies that there exist kinematically
well-defined questions about future continuation that cannot be resolved within
any admissible encoding. This observation underlies the computational and
predictive limits analyzed in a subsequent paper.
\end{remark}

% ============================================================
% Next section will be:
% \section{Summary and Relation to Downstream Encodings}
% ============================================================
\section{Summary and Relation to Downstream Encodings}

In this paper we have developed a notion of kinematics intrinsic to the encoding
atlas of Modal Triplet Theory. Without assuming spacetime, geometry, or
equations of motion, we have shown how configuration, position, motion, causal
structure, horizons, selection, and irreversibility arise as structural features
of admissible continuation across overlapping reduced descriptions.

The central kinematic insight is that persistence replaces displacement as the
primitive notion of motion. Coherent structures move insofar as their
representation persists across chains of admissible encodings. Worldlines are
equivalence classes of such continuation chains, and causal structure is induced
by the partial ordering defined by admissibility. Null and timelike distinctions,
horizons, and termination of motion follow directly from the structure of
admissible continuation and finite admissibility.

These results are kinematic rather than dynamical. No metric, connection, or
field equation has been introduced. Instead, kinematic structure is shown to be
a consequence of the encoding atlas itself. Geometry and dynamics, when they
appear in later papers, are not taken as defining kinematics but as encoding
responses required to stabilize and organize the kinematic features identified
here.

In particular:
\begin{itemize}
\item the stabilization of causal cones and null boundaries motivates the
emergence of geometric and gravitational encoding structures;
\item the organization of redundancy in admissible continuation motivates gauge
encoding structures;
\item the presence of selection fronts and horizons motivates discrete and
quantized encoding constraints.
\end{itemize}

The present paper therefore occupies a precise position in the Modal Triplet
Theory series. It depends only on the structural core and provides the kinematic
interpretation necessary for understanding gravity, gauge structure,
quantization, and unified encodings as downstream responses to structural
obstructions rather than as fundamental postulates.

Together with the structural core and the obstruction classification developed
elsewhere, this kinematic framework completes the pre-geometric layer of the
theory.

\begin{remark}[Series progression]
The logical progression of the series is as follows:
\begin{enumerate}[label=(\roman*)]
\item the structural core establishes admissibility and the impossibility of
global description;
\item coherent kinematics defines motion, causality, and irreversibility without
spacetime;
\item obstruction classification explains why circle, lens, and nil are
inevitable;
\item encoding-class papers develop gravity, gauge structure, and quantization
as bookkeeping responses;
\item realization papers instantiate these encodings geometrically.
\end{enumerate}
\end{remark}

% ============================================================
% End of Paper A1

\end{document}
